% 05-signatures.tex: криптографические свидетельства.
\section{Удостоверение истины: подписи, хранение ключей, два свидетеля и одна миграция}
\label{sec:signatures}

\motif{Запись в реестре не есть доказательство того, кто её сделал.}

\analogy{Каждая запись в реестре несёт печать сделавшего её чиновника, оттиснутую штемпелем, которого больше никто не держит. Два пробирщика, выученные в разных школах, осматривают каждую новую печать, и печать, которую признаёт лишь один из них, не принимается.}

\subsection{Для чего нужна подпись}

Цифровая подпись есть печать, поставить которую способен лишь владелец закрытого ключа и проверить
которую способен всякий, у кого есть парный открытый ключ. Полярис подписывает дайджест ША3-256
значения токена ключом МЛ-ДСА выдающего ведомства и хранит подпись рядом с открытым ключом,
который её поставил. Утверждение, которое несёт подпись, узко и постоянно: \emph{этот ключ
подписал эти байты.} Оно ничего не говорит о том, сохраняются ли у токена полномочия, что есть
задача следующего раздела, и ничего о том, кто владеет ключом, что есть задача хранения. \sImpl

Алгоритм по умолчанию есть МЛ-ДСА-65 (ФИПС 204, категория стойкости 3), а рядом с ним МЛ-ДСА-87
(категория 5); МЛ-ДСА-44 отвергается повсюду как лежащий ниже предела. Ни одно подписанное тело не
зашивает в себя свой алгоритм: хранимый ключ решает это перед подписанием, потому что
\code{алгоритм} сам есть подписываемое поле, а реестр объявляет принимаемый набор. СЛХ-ДСА-128с и
СЛХ-ДСА-256с, хеш-основанные запасные варианты, не опирающиеся на решёточные допущения,
зарегистрированы в таблице алгоритмов именно потому, что опираются на одни лишь хеш-функции, так
что поворот прочь от решёток есть обновление строки; ни один подписант СЛХ-ДСА не подключён, и
книга учёта состояния несёт этот пробел. \sImpl{} для МЛ-ДСА; \sFut{} для хеш-основанного
подписанта.

\subsection{Хранение ключей}

Закрытый ключ никогда не лежит в коде приложения. Он лежит за интерфейсом хранения с тремя
драйверами: файл (разработка), ПККС\#11 (аппаратный модуль безопасности или программный токен;
ключ вырабатывается внутри токена и наружу не выходит) и облачная служба управления ключами.
Непрерывная интеграция прогоняет драйвер ПККС\#11 против программного токена, реализующего МЛ-ДСА,
а не против аппаратного модуля, и проворачивает внутри него замену ключа: токен, подписанный
старым внутритокенным ключом, по-прежнему проверяется после переключения, а новый ключ уже
подписывает. Конвертная криптография драйвера СУК настоящая, а его служба есть заместитель. Один
промышленный профиль делает модуль единственным подписантом: приложение отказывается
запускаться, если присутствует файловый ключ или настоящее подписание недоступно. Какой драйвер,
какой модуль и кто держит ключ, суть решения, которые книга готовности возлагает на
развёртывающую организацию. \sImpl{} для драйверов и учения; \sExt{} для аппаратного хранения,
поскольку ни один аппаратный модуль не применялся.

Заглушка заслуживает отдельного абзаца, потому что именно здесь честную систему легче всего
прочесть неверно. Без флага настоящего
подписания и без установленной решёточной библиотеки подписывающий модуль производит
детерминированную привязку ША3-256 с пометкой
\code{ДЕТЕРМИНИРОВАННАЯ-ЗАГЛУШКА-ША3-256}: тридцать два байта, не проверяемые ни против какого
ключа, тогда как настоящая подпись МЛ-ДСА-65 занимает 3309. Это и есть значение по умолчанию при
разработке и на большей части непрерывной интеграции. Оно ограждено: промышленный запуск
падает закрытым, если настоящее подписание действительно недоступно, заглушка
есть именованный профиль разработки, и она громко предупреждает, когда её применяют безымянно; и
всё же именно из-за неё репозиторий больше не называет себя постквантовым, не оговорив этого тут
же. \sImpl

\subsection{Два свидетеля при выдаче, один при использовании с выборочной проверкой и один внешний набор}

Принцип двух свидетелей гласит, что всякий криптографический вердикт, который Полярис выпускает
наружу, обязан быть проверяем второй, независимой реализацией, написанной отдельно, представленной
иначе, не разделяющей ни строки кода и либо соглашающейся с вердиктом, либо явно воздерживающейся.
Верификатор, согласный сам с собой, не установил ничего. При выдаче только что произведённая
подпись проверяется двумя независимыми реализациями ФИПС 204, решёточной библиотекой и
библиотекой на основе ОпенЭсЭсЭл, и ей отказывают в сохранении, если её не принимают обе;
отсутствующий свидетель при выдаче есть отказ, а не понижение. \sImpl

Проверка при использовании прогоняет одного свидетеля, потому что про сохранённую подпись уже
известно, что она проверяется обоими, а один свидетель обнаруживает подмену примерно вдесятеро
быстрее: на эталонной машине около 7950 проверок в секунду на ядро против примерно 740 при двух
свидетелях. Быстрый путь зарабатывает свою скорость тем, что его непрерывно проверяют:
обязательная случайная доля успешных проверок переигрывается вторым свидетелем, расхождение
увеличивает счётчик и вызывает пейджер с наивысшей важностью, а каждый ответ называет, какой набор
свидетелей отработал. \sImpl

Одна вещь про свидетелей не замкнута на себя. При каждой отправке векторы проверки МЛ-ДСА-65 из
проекта Вайчпруф, сторонний набор, привязанный к коммиту, в котором принятие подписи, названной
набором недействительной, есть отказ, прогоняются под обоими свидетелями: 58 из 58 совпадают. Это
свидетельство того, что проверка примитива в Полярис сходится с внешним авторитетом на 58 случаях.
Это не кошелёк, не результат по совместимости, и оно ничего не говорит о протоколе удостоверений;
сводная таблица репозитория записывает это ровно этими словами. \sWit{} для примитива.

Тот же принцип управляет вердиктом с нулевым разглашением (раздел~\ref{sec:privacy}): повторная
реализация на Питон с обыкновенными целыми числами по модулю простого, не разделяющая ни строки
кода с доказывателем на Раст, заново выводит утверждение о принадлежности и обязана согласиться, а
воздерживается она открыто по той единственной оси, которую смоделировать не в силах, по
целостности байтов доказательства. Что принцип ловит, так это расхождение реализаций. Чего он не
ловит, так это общего неверного прочтения спецификации, и репозиторий говорит об этом рядом с
каждым утверждением о двух свидетелях. Свидетели также проверяются мутацией,
так что утверждение, что две реализации обязаны согласиться, есть утверждение и о тестах, а не
только о коде.

\subsection{Каноническое подписание: одни и те же байты с обеих сторон}

Всякий подписанный предмет в Полярис есть ДЖСОН-утверждение с упорядоченными ключами и без
пробельных символов, взятое дайджестом ША3-256 и затем подписанное. Приложение строит
утверждение; отделяемый верификатор (раздел~\ref{sec:verifiers}) перестраивает его из предмета,
чтобы проверить подпись. Два построения обязаны совпадать побайтово, иначе подлинный предмет
проверку не пройдёт, и тест, оракул канонической равнозначности, закрепляет каждый подписываемый
тип с обеих сторон. Единственная ошибка, ради которой оракул и существует, была найдена на ревью:
предмет, поле \code{алгоритм} которого выставлялось после подписания, нёс подписываемое поле, не
покрытое подписью. \sImpl

\subsection{У ключей есть жизненный цикл}

Реестр ключей есть таблица событий, работающая только на добавление: зарегистрирован, выведен из
обращения или скомпрометирован с момента, который может предшествовать обнаружению, и из неё
публикуется подписанный список доверия. Верификатор, держащий список доверия, определяет состояние
ключа на любой момент, так что удостоверение под скомпрометированным ключом эмитента отвергается
независимо от собственной описи эмитента, которую вор, завладевший ключом, мог бы подделать, а
документ, подписанный и помеченный временем до компрометации, остаётся действительным, тогда как
подписанный после не остаётся. Учение на двух экземплярах записывает компрометацию в одном органе
и наблюдает, как решение другого переворачивается по проводу, без всякого содействия
скомпрометированной стороны. \sImpl

\subsection{Гибкость при проверяемом пути миграции и во что обошёлся бы взлом}

Стойкость МЛ-ДСА-65 не опирается на предсказание о квантовом железе; она опирается на трудность Модуль-ОСО и
Модуль-КЦЗ, решёточных задач, на которых построена схема, против лучших известных классических и квантовых алгоритмов, а это может оказаться
неверным по математическим причинам. Два настоящих риска суть криптоаналитический прорыв против
решёточных допущений и изъян реализации, и устройство уже отвечает на оба, хеш-основанным
подстрахованием в реестре и двумя свидетелями при выдаче. Репозиторий способен показать, а потому и утверждает, \emph{алгоритмическую гибкость при
проверяемом пути миграции}: С7 делает алгоритм полем, \entity{ДопускКАлгоритму} разрешает его
для каждого органа, а население можно переподписать под новым набором параметров путём,
выполняемым при каждой отправке. \sImpl

Этот путь измерен, а не запланирован. \code{полярис мигрировать-население} переподписывает всё
действующее население под новым набором параметров и построен из тех же ограничений, на которые
опирается потокенная процедура: уникальное ограничение есть точка упорядочивания, так что два
исполнителя, состязающиеся за одно удостоверение, производят одну строку, а шестьдесят четыре
исполнителя делят население без всякого согласования. Значение, которое важно, есть не
пропускная способность, а то, сколько держателей держат удостоверение, не проверяемое ничем, и
учение берёт эту выборку на каждой границе пакета, а не по концам, потому что разрыв, который
открывается и закрывается между двумя конечными точками, невидим для проверки <<до и после>>. Он
равен нулю на всём протяжении. Старая подпись продолжает проверяться до своей даты устаревания, и
этот промежуток \emph{и есть} миграция: объявить вытесненный алгоритм устаревшим отказываются,
пока хоть одно действующее удостоверение не мигрировало, потому что иначе держатель обнаружил бы
разрыв на границе, пока консоль сообщала бы об успехах. Под настоящим МЛ-ДСА-87 с одним хранимым
ключом учение переподписывает около 345 удостоверений в секунду на одном исполнителе, и 91 процент
этого времени занимает подписание, так что население в 350 миллионов есть около двенадцати дней на
одном исполнителе и четыре с половиной часа на шестидесяти четырёх; закупка мощности базы данных
под это почти ничего не даёт. \sImpl

Во что обошёлся бы взлом Модуль-ОСО, записано в книге готовности, и настоящий документ это
повторяет. Он не потребовал бы изменения схемы, изменения на пути проверки или новой модели
доверия. Он стоил бы всякого удостоверения, уже выданного под взломанным алгоритмом, и ещё одной
вещи, которой никакая гибкость не возвращает: во время переключения удостоверение несёт
классическую подпись рядом с постквантовой, а классическую половину не защищает ничто из того, что
даёт это устройство, так что удостоверение, собранное сегодня, станет читаемым в день падения его
классической подписи, на что бы система ни мигрировала потом. Эта несимметричность и есть причина,
по которой подписывать постквантово стоит уже сейчас, и она же есть предел того, чего этим
достигают.

Собственная лаборатория репозитория затем спросила, на той ли стороне находится гибкость, которой
она нужна, и обнаружила, что гибкость есть только со стороны эмитента. Поле
\code{дата\_устаревания} алгоритма не доходит ни до одного подписанного предмета: реестр публикует
голый перечень имён алгоритмов, список доверия несёт состояние только по каждому ключу, а
принимаемый набор отделяемого верификатора есть словарь в коде. Поэтому верификатор не способен
узнать, что алгоритм объявлен устаревшим, вывод алгоритма из обращения есть выпуск программного
обновления для каждой доверяющей стороны, а не миграция, и отзыв ключа этого не заменяет, потому
что если допущение рухнет, нападающий подделает подпись под любым ключом, который список доверия
называет действующим. Откат и смешанное окно во время частичной миграции не измерены. Изменений не
предлагается; находка записана как ограничение для развёртывающей организации. \sExt

\subsection{Что остаётся классическим}

Подпись токена, дайджесты, обязательства Меркла, доказательство с нулевым разглашением, обмен
ключами от клиента до кромки для современных клиентов и переход от
приложения к пулеру соединений суть постквантовые, причём последние два считаны с настоящего
рукопожатия, а не выведены из базового образа. Переход от пулера к базе данных, подписи
сертификатов и учётные данные ВебАутн оператора (вход по аппаратному ключу) остаются классическими, каждое из этого упирается
в третью сторону (выпуск ПостгреСКЛ, публичная экосистема сертификатов, железо аутентификаторов) и
каждое изложено в книге учёта состояния со своей угрозой и своей открытостью. Ни решёточная
библиотека, ни система доказательств, ни схема не проходили внешнего криптографического аудита.
\sExt

\begin{layercard}{Карточка слоя: криптографические свидетельства}
\qa{Что это}{Подписи МЛ-ДСА над дайджестами ША3-256, хранимые вместе со своими открытыми ключами; интерфейс хранения с драйверами файла, ПККС\#11 и СУК; две независимые проверяющие реализации и один внешний набор векторов; канонический формат утверждения; реестр ключей только на добавление и подписанный список доверия; работающая миграция населения.}
\qa{Зачем оно существует}{Утверждению, покидающему базу данных, нужно везти доказательство с собой; доказательство, проверить которое способна лишь одна программа, есть обещание; а алгоритм, который нельзя заменить, есть назначенный срок.}
\qa{Чему доверяет}{ФИПС 204 и 202 в том виде, как они стандартизованы; тому, что две библиотеки не разделяют одной ошибки; той подсистеме хранения, которую выбрал оператор.}
\qa{Чему отказывает}{Вердикту одной реализации при выдаче; алгоритму, названному в коде, а не ключом; подписи, сохранённой без своего ключа; заглушке, выдающей себя за подпись; объявлению устаревания, пока хоть одно удостоверение осталось бы не проверяемым ничем.}
\qa{Что видит}{Байты, которые подписывает, а они суть дайджесты: путь подписания никогда не видит человека.}
\qa{Чего не видит}{Сохраняются ли у подписанного утверждения полномочия (раздел~\ref{sec:authorization}); знает ли верификатор в поле, что алгоритм выведен из обращения.}
\qa{Если откажет}{Поддельная подпись падает у обоих свидетелей; скомпрометированный ключ выводится из обращения с указанного момента через реестр, и всякий верификатор, держащий список доверия, отвергает подписанное им после; взломанный алгоритм стоит каждого удостоверения под ним и переживается переподписанием, со измеренной скоростью.}
\qa{Сейчас или потом}{Сейчас. Аппаратный модуль и внешний аудит потом, и принадлежат они развёртывающей организации; подписанное состояние алгоритма, которое мог бы прочесть верификатор, не построено.}
\qa{Внутри и снаружи}{Внутри: схема, хранящая строки подписей. Снаружи: независимые верификаторы, проверяющие их без эмитента.}
\end{layercard}

\nextlayer{Подпись, которую проверяет собственный сервер эмитента, не доказывает постороннему
ничего, потому что постороннему приходится доверять серверу. Следующий слой забирает проверку из
рук эмитента целиком и, впервые, отдаёт её в руки программ, написанных не автором эмитента.}
